%------------------------------
% Packages and Themes
%------------------------------

\documentclass[aspectratio=169]{beamer}
% faculty: medicine, engineering, math, science, or default
% progressbar: opt-in progress bar under each frame title (omit to disable)
\usetheme[faculty=engineering,progressbar]{SimpleUUlm}

% Packages
\usepackage{graphicx}
\usepackage{booktabs}
\usepackage{tabularx}
\usepackage{hyperref}

% Custom column type and commands
\newcolumntype{L}{>{\raggedright\arraybackslash}X}
\newcommand{\swatch}[1]{%
    \fcolorbox{UUlmMutedBlack}{#1}{\rule{0pt}{1.4em}\hspace{2.2em}}%
    \\[2pt]{\scriptsize\texttt{#1}}}

%------------------------------
% Title Page
%------------------------------

\title{SimpleUUlm Beamer Theme}
\subtitle{A simple and clean beamer theme for the University of Ulm}

\author{Moritz Wieland}

% Choose from:
% Faculty of Engineering, Computer Science and Psychology
% Faculty of Mathematics and Economics
% Faculty of Natural Sciences
% Faculty of Medicine
\institute{
    Faculty of Engineering, Computer Science and Psychology \\
    Ulm University
}

\date{\today}

%------------------------------
% Presentation Slides
%------------------------------

\begin{document}

\begin{frame}[noframenumbering, plain]
    \titlepage
\end{frame}

\begin{frame}{Overview}
    \tableofcontents
\end{frame}

%--------------------------------------------------
\section{Text, Lists and Emphasis}
%--------------------------------------------------

\begin{frame}{Bullet Lists}
    Bullets and nested lists pick up the accent color.

    \begin{itemize}
        \item Top-level item with a nested list
        \begin{itemize}
            \item Second-level item
            \item Another second-level item
        \end{itemize}
        \item Inline math renders too: $e^{i\pi} + 1 = 0$
    \end{itemize}
\end{frame}

\begin{frame}{Numbered Lists}
    Enumerations use the same accent color for their markers.

    \begin{enumerate}
        \item First step
        \item Second step
        \item Third step
    \end{enumerate}
\end{frame}

\begin{frame}{Description Lists}
    Description lists are themed too --- the term picks up the accent color.

    \begin{description}
        \item[Accent] Driven by the selected faculty.
        \item[Structure] Titles, bullets, and rules share one color.
        \item[Fira Sans] The theme ships with the FiraSans font family.
    \end{description}
\end{frame}

\begin{frame}{Emphasis and Alerts}
    The accent color follows the faculty chosen via
    \texttt{\textbackslash usetheme[faculty=...]}.
    Use \alert{alert text} for warnings and
    \textcolor{UUlmInstitute}{accent text} for emphasis.

    \vspace{0.5em}
    Inline styling: \textbf{bold}, \emph{emphasis}, \texttt{monospace}, and
    a \href{https://github.com/mowi12/SimpleUUlm-BeamerTheme}{hyperlink}.
\end{frame}

\begin{frame}{Blocks}{Standard, alert, and example variants}
    \begin{block}{Standard block}
        Title bar in the accent color, body in a light tint of it.
    \end{block}

    \begin{alertblock}{Alert block}
        Always red, independent of the faculty color --- for critical points.
    \end{alertblock}

    \begin{exampleblock}{Example block}
        Sand accent --- for analogies, asides, or worked examples.
    \end{exampleblock}
\end{frame}

%--------------------------------------------------
\section{Color Palette}
%--------------------------------------------------

\begin{frame}{Named Colors}{The full UUlm palette at a glance}
    Every color is available by name; \texttt{UUlmInstitute} is an alias that
    tracks the chosen faculty.

    \vspace{1em}
    \begin{columns}[t]
        \column{0.24\textwidth}\centering
            \swatch{UUlmRed}\\[1em]
            \swatch{UUlmGreen}
        \column{0.24\textwidth}\centering
            \swatch{UUlmDarkBlue}\\[1em]
            \swatch{UUlmOrange}
        \column{0.24\textwidth}\centering
            \swatch{UUlmLightBlue}\\[1em]
            \swatch{UUlmAccent}
        \column{0.24\textwidth}\centering
            \swatch{UUlmMutedBlack}\\[1em]
            \swatch{UUlmInstitute}
    \end{columns}
\end{frame}

%--------------------------------------------------
\section{Columns}
%--------------------------------------------------

\begin{frame}{Two Columns}
    \begin{columns}[t]
        \begin{column}{0.48\textwidth}
            \textbf{Left column}
            \begin{itemize}
                \item Columns split the frame
                \item Good for text beside a figure
                \item Or for side-by-side comparison
            \end{itemize}
        \end{column}
        \begin{column}{0.48\textwidth}
            \begin{block}{Right column}
                A block sits comfortably inside a column for highlighting a
                takeaway next to supporting text.
            \end{block}
        \end{column}
    \end{columns}
\end{frame}

%--------------------------------------------------
\section{Tables and Figures}
%--------------------------------------------------

\begin{frame}{Full-width Table}
    \small
    \begin{tabularx}{\textwidth}{@{} l L L @{}}
        \toprule
        \textbf{Feature} & \textbf{Description} & \textbf{Notes} \\
        \midrule
        Frame title rule & Rule under each frame title & Optional subtitle above it \\[4pt]
        Section page     & \texttt{\textbackslash AtBeginSection} divider in this deck & White on accent \\[4pt]
        Frame number     & Plain number in the footer & From \texttt{footline[frame number]} \\[4pt]
        Faculty color    & One command swaps the whole accent & Five presets \\
        \bottomrule
    \end{tabularx}

    \vspace{2em}
    {\small\color{UUlmMutedBlack}%
        \texttt{L} is a left-aligned, wrapping tabularx column from the preamble
    }
\end{frame}

\begin{frame}{Comparison Table}
    \small
    \begin{tabularx}{\textwidth}{@{} l >{\centering\arraybackslash}X
                                        >{\centering\arraybackslash}X @{}}
        \toprule
                  & \textbf{Option A} & \textbf{Option B} \\
        \midrule
        Setup     & \textcolor{UUlmInstitute}{Simple}
                  & \textcolor{UUlmRed}{Involved} \\[3pt]
        Control   & \textcolor{UUlmRed}{Limited}
                  & \textcolor{UUlmInstitute}{Full} \\[3pt]
        Verdict   & Good default
                  & For advanced use \\
        \bottomrule
    \end{tabularx}

    \vspace{2em}
    {\small\color{UUlmMutedBlack}%
        Accent and red can be mixed inline to encode good / bad.%
    }
\end{frame}

\begin{frame}{Figure with Caption}
  \begin{figure}
    \rule{3cm}{2cm}              % stand-in for an image
    \caption{The label before this text uses the accent caption color.}
  \end{figure}
\end{frame}

%--------------------------------------------------
\section{Links, Buttons and Notes}
%--------------------------------------------------

\begin{frame}{Buttons}\label{buttons}
  \beamerbutton{A button}\hfill
  \hyperlink{buttons}{\beamergotobutton{Self link}} % same styling, with arrow

  \vspace{1em}
  Buttons use the accent background with a matching border.
\end{frame}

\begin{frame}{Footnote}
  Some claim.\footnote{The mark up there uses the accent footnote-mark color.}
\end{frame}

%--------------------------------------------------
\section{Wrap-up}
%--------------------------------------------------

\begin{frame}{Conclusion}
    \begin{itemize}
        \item One theme, five faculty presets via \texttt{faculty=...}
        \item Themed titles, lists, blocks, captions, buttons, and footnotes
        \item Tables, columns, and the FiraSans font ready out of the box
    \end{itemize}

    \vspace{0.6em}
    \begin{alertblock}{}
        Swap \texttt{\textbackslash usetheme[faculty=engineering]} for
        \texttt{medicine}, \texttt{math}, \texttt{science}, or
        \texttt{default} to retheme the whole deck.
    \end{alertblock}
\end{frame}

%--------------------------------------------------
\subsection{References}
%--------------------------------------------------

\begin{frame}{Citing Sources}
    Cite references inline with \texttt{\textbackslash cite}: as shown by
    Smith~\cite{r1}, citations link to the bibliography below.
\end{frame}

\begin{frame}[allowframebreaks]{References}
    \footnotesize
    \bibliographystyle{apalike}
    \bibliography{reference}
\end{frame}

\end{document}
